%----------
%	EVALUACIÓN E INSTRUMENTACIÓN
%----------

\subsection{Métricas}

Las predicciones se comparan con la anotación experta sobre las 1.313 piezas
etiquetadas. Cada variable se trata como un problema binario en el que la clase
positiva es la presencia del fenómeno, es decir, el \textit{Sí}. De esa
comparación resultan cuatro recuentos: los verdaderos positivos ($VP$), piezas que
tanto el sistema como la anotadora marcan, los falsos positivos ($FP$), que el
sistema marca y la anotadora no, los falsos negativos ($FN$), que la anotadora
marca y el sistema pasa por alto, y los verdaderos negativos ($VN$), en que ambas
coinciden en la ausencia. Sobre esos cuatro valores se construyen todas las
métricas empleadas.

\paragraph{Exactitud.} Es la proporción de piezas en que sistema y anotadora
coinciden, con independencia de la clase:
\begin{equation}
    \text{Exactitud} = \frac{VP + VN}{VP + VN + FP + FN}
\end{equation}
Su ventaja es que resulta inmediatamente interpretable, ya que responde a la
pregunta de cuántas veces acierta el sistema. Su inconveniente es que la domina la
clase mayoritaria. En V33, con un 6,2\,\% de positivos, no marcar nada produce una
exactitud del 93,8\,\% sin haber detectado un solo caso
(Sección~\ref{sec:iris_variables}), de manera que por sí sola no distingue un
sistema competente de uno que no lo es.

\paragraph{Precisión.} De todo lo que el sistema marca, mide qué proporción estaba
efectivamente marcada por la anotadora:
\begin{equation}
    \text{Precisión} = \frac{VP}{VP + FP}
\end{equation}
Responde a la pregunta de cuánto cabe fiarse de una alerta del sistema, lo que la
hace directamente relevante para una herramienta de asistencia, ya que una
precisión baja se traduce en avisos que la persona analista deberá descartar. Su
inconveniente es que se calcula solo sobre lo marcado, así que un sistema muy
conservador puede alcanzar una precisión alta con un puñado de aciertos, y
queda indefinida si no marca nada.

\paragraph{Recall.} También llamado sensibilidad, mide qué proporción de lo que la
anotadora marcó consigue detectar el sistema:
\begin{equation}
    \text{Recall} = \frac{VP}{VP + FN}
\end{equation}
Responde a la pregunta complementaria, la de cuánto sexismo se está dejando sin
señalar. Su inconveniente es simétrico al de la precisión, puesto que marcarlo
todo garantiza un recall de 1 sin discriminar nada. Precisión y recall solo tienen
sentido, por tanto, leídos conjuntamente.

\paragraph{F1.} Resume ambas en su media armónica, que penaliza el desequilibrio
entre las dos más de lo que lo haría una media aritmética:
\begin{equation}
    F_1 = 2 \cdot \frac{\text{Precisión} \cdot \text{Recall}}{\text{Precisión} + \text{Recall}}
\end{equation}
La variante \textbf{macro} calcula el $F_1$ de cada clase por separado y promedia
los dos valores sin ponderar por su frecuencia, de modo que la clase minoritaria
pesa tanto como la mayoritaria. Es una corrección importante frente a la
exactitud, pero no elimina el problema del desequilibrio, ya que en una variable
con un 81\,\% de positivos como V25 marcarlo todo produce un $F_1$ de 0,895 sobre
la clase positiva.

\paragraph{Coeficiente $\kappa$ de Cohen.} Mide el acuerdo entre dos anotadores
descontando el que cabría esperar por azar~\cite{cohen1960}:
\begin{equation}
    \kappa = \frac{p_o - p_e}{1 - p_e}
\end{equation}
donde $p_o$ es el acuerdo observado, equivalente a la exactitud, y $p_e$ el
acuerdo esperado si ambos etiquetasen al azar respetando sus respectivas
frecuencias marginales. Un valor de 1 indica acuerdo perfecto, un 0 indica un
acuerdo equivalente al del azar y los valores negativos señalan un desacuerdo
sistemático. Su ventaja es que sitúa el resultado frente a una referencia
razonable en lugar de frente al cero absoluto, y que procede del análisis de
contenido, disciplina de la que nace este trabajo. Su inconveniente es la conocida
\textbf{paradoja de $\kappa$}~\cite{feinstein1990}, según la cual con
distribuciones muy desequilibradas el acuerdo esperado por azar se dispara y el
coeficiente se desploma aunque la exactitud sea alta. Esto vuelve los valores de
$\kappa$ difícilmente comparables entre variables de prevalencia distinta.

Ninguna métrica basta por sí sola, y la que resulta informativa depende de la
variable. En V33 y V35, donde los positivos son escasos, la exactitud carece de
valor, la precisión se vuelve inestable por apoyarse en muy pocos casos y el
recall es el indicador decisivo, ya que revela cuánto se queda sin detectar. En
V25 y V26, con prevalencias del 81\,\%, ocurre lo contrario, pues el recall y el
$F_1$ se inflan con facilidad y lo informativo pasa a ser la capacidad de
reconocer las piezas sin sexismo, que es donde $\kappa$ resulta más exigente. Por
este motivo los resultados del Capítulo~\ref{results_and_discussion} se presentan
siempre acompañados de la prevalencia de cada variable y del rendimiento del
clasificador trivial (Sección~\ref{sec:iris_variables}), que fijan la referencia
mínima frente a la que interpretarlos.

Existe, además, un criterio que ordena esta elección y que procede del uso
previsto del sistema. Como herramienta de apoyo a la postedición, su cometido
consiste en señalar a la persona analista las piezas que merecen una segunda
lectura, de modo que un falso positivo cuesta unos segundos de descarte mientras
que un falso negativo supone un caso de sexismo que nadie revisará. Esa asimetría
otorga al recall una prioridad que no tendría en un clasificador autónomo, y así se
discute en el Capítulo~\ref{results_and_discussion}.


\subsection{Evaluación \textit{zero-shot} sobre el corpus completo}

La evaluación se realiza en régimen \textit{zero-shot}
(Sección~\ref{sec:zeroshot_iris}), lo que significa que el modelo recibe la
metodología en el \textit{prompt} y que no se ajusta ningún parámetro a la vista
de los datos. Dos consecuencias de ese régimen, ya justificadas en la
Sección~\ref{sec:zeroshot_iris}, condicionan la lectura de los resultados. En
cuanto a la adquisición de conocimiento, el sistema no incorpora información del corpus, de modo que su competencia procede íntegramente del preentrenamiento y de la
metodología que se le entrega en cada llamada. En cuanto al modelado estadístico,
no estima probabilidades a partir de los datos anotados ni ajusta su umbral a la
prevalencia de cada variable, así que su propensión a marcar refleja su
interpretación del criterio y no la distribución real de las clases.

De ahí se desprende una consecuencia práctica sobre el diseño de la evaluación.
Todas las cifras del Capítulo~\ref{results_and_discussion} se calculan sobre las
1.313 piezas etiquetadas, sin reservar un subconjunto de desarrollo o de prueba.
Al no existir entrenamiento ni selección de umbrales sobre los datos, no cabe el
riesgo de sobreajustar a un subconjunto y medir después el rendimiento en él, y el
corpus completo puede emplearse como material de evaluación.


\subsection{Ejecución a escala: enrutado y granja}

La misma secuencia de procesos sirve para todos los modelos gracias al enrutamiento por
identificador (Sección~\ref{sec:llms_models_iris}). Para los modelos comerciales,
las peticiones van a la API correspondiente. Para los locales, se ejecutan sobre
la granja de GPU del Departamento de Teoría de la Señal y Comunicaciones mediante
\textit{Ollama}, repartiendo el corpus en fragmentos (\textit{shards}) que se procesan
en paralelo sobre varios servidores, con las peticiones lanzadas desde un equipo
cliente del laboratorio. El coste computacional de esta vía local se cuantifica en
la Sección~\ref{subsec:iris_coste}. El código se recoge en el
Apéndice~\ref{appendix_code}.
